\subsection{Noetherian topological spaces}

DEFINITION 3. A topological space X is called Noetherian if every non-empty set of closed subsets of X, ordered by inclusion, has a minimal element.

It amounts to the same to say that every non-empty set of open subsets of X, ordered by inclusion, has a maximal element, or that every decreasing (resp. increasing) sequence of closed (resp. open) sets is stationary (Set Theory, Chapter III, § 6, no. 5, Proposition 6).

PROPOSITION 8. (i) Every subspace of a Noetherian space is Noetherian.

\begin{enumerate}
\def\labelenumi{(\roman{enumi})}
\setcounter{enumi}{1}
\item
  Let \((\mathbf{A}_i)_{i\in I}\) be a finite covering of a topological space \(\mathbf{X}\) . If the subspaces \(\mathbf{A}_i\) of \(\mathbf{X}\) are Noetherian, \(\mathbf{X}\) is Noetherian.

\end{enumerate}

(i) Let X be a Noetherian space, A a subspace of X and \((\mathbf{F}_{n})\) a decreasing sequence of subsets of A which are closed in A; then \(F_{n} = \overline{F}_{n} \cap A\) and the closures \(\overline{F}_{n}\) of the \(F_{n}\) in X form a decreasing sequence of closed subsets of X. As this sequence is stationary, so is the sequence \((\mathbf{F}_{n})\) .

(ii) Let \((\mathbf{G}_n)_{n\geqslant 0}\) be a decreasing sequence of closed subsets of \(\mathbf{X}\) ; by hypothesis, each of the sequences \((\mathrm{G}_n\cap \mathrm{A}_i)_{n\geqslant 0}\) is stationary. As I is finite, there is an integer \(n_0\) such that, for \(n\geqslant n_0\) , \(\mathbf{G}_n\cap \mathbf{A}_i = \mathbf{G}_{n_0}\cap \mathbf{A}_i\) for all \(i\in \mathbf{I}\) . But

\[
\mathbf {G} _ {n} = \bigcup_ {i \in I} \left(\mathbf {G} _ {n} \cap \mathbf {A} _ {i}\right)
\]

and hence the sequence \((\mathbf{G}_n)\) is stationary and \(\mathbf{X}\) is Noetherian.

PROPOSITION 9. For a topological space X to be Noetherian, it is necessary and sufficient that every open set in X be quasi-compact.

To show that the condition is necessary, it is sufficient, by virtue of Proposition 8, to prove that every Noetherian space X is quasi-compact. Let \((\mathbf{U}_{\iota})_{\iota\in\mathbf{I}}\) be an open covering of X; the set of finite unions of sets \(U_{\iota}\) is non-empty and hence admits a maximal element \(V=\bigcup_{\iota\in H}U_{\iota}\) , where H is a finite subset of I. By definition, \(V\cup U_{\iota}=V\) for all \(\iota\in I\) and hence V=X.

Conversely, suppose that every open set of X is quasi-compact and let \((\mathbf{U}_{n})\) be an increasing sequence of open subsets of X. The union V of the \(U_{n}\) is open and hence quasi-compact; as \((\mathbf{U}_{n})\) is an open covering of V, there is a finite subfamily of \((\mathbf{U}_{n})\) which is a covering of V and hence \(V = U_{n}\) for some index n, which proves that the sequence \((\mathbf{U}_{n})\) is stationary.

LEMMA 1 (``Principle of Noetherian induction''). Let E be an ordered set every non-empty subset of which admits a minimal element. Let F be a subset of E with the following property: if \(a \in E\) is such that the relation \(x < a\) implies \(x \in F\) , then \(a \in F\) . Then \(F = E\) .

Suppose \(\mathbf{F} \neq \mathbf{E}\) ; then \(\complement\mathbf{F}\) would have a minimal element \(b\) . By definition, \(x \in \mathbf{F}\) for all \(x < b\) , which implies that \(b \in \mathbf{F}\) , which is a contradiction.

PROPOSITION 10. If X is a Noetherian space, the set of irreducible components of X (and a fortiori the set of connected components of X) is finite.

It is sufficient to prove that X is a finite union of irreducible closed subsets (no. 1, Proposition 6). Let us show that the principle of Noetherian induction can be applied taking E to be the set of closed subsets of X, ordered by inclusion, and F to be the set of finite unions of irreducible closed subsets. Let Y be a closed subset of X such that every closed subset \(\neq Y\) of Y belongs to F. If Y is irreducible, then \(Y \in F\) by definition; otherwise, Y is the union of two closed subsets \(Y_{1}, Y_{2}\) which are distinct from Y. Then \(Y_{1} \in F\) and \(Y_{2} \in F\) by hypothesis, whence \(Y \in F\) by definition of F.

In particular it follows that a Hausdorff Noetherian space is necessarily finite.

